\section{Methods: Group 4 --- Sagittal Alignment}
\label{sec:g4}

Group 4 computes the global cervical Cobb angle (C2--C7 / C3--C7), per-level segmental angles, a
posterior-tangent cross-check, and a lordosis classification. The Cobb angle is the most technically
demanding measurement in the system, and its method is the project's main methodological story.

\paragraph{The corner problem.} A Cobb angle is the angle between the top and bottom endplate tangents.
Deriving those tangents from four/six corner landmarks is unstable on lordotic cervical necks because
endplates are concave and sloped~\cite{chen2013}: the legacy corner method read healthy necks as
\SI{-21}{\degree} (kyphotic) when they are in fact lordotic. The literature is unambiguous that fitting a
line to the full endplate boundary beats corner extrema~\cite{wang2023}.

\paragraph{Three methods, evaluated.} We implemented and compared three Cobb methods on the 12 healthy
necks (precision/coverage, since no GT exists):
\begin{enumerate}[leftmargin=*]
  \item \textbf{Canal-cut endplate-line} --- fit the tangent to the canal-cut body's endplate boundary.
  Fixes the \emph{sign} (now lordotic) but the C6/C7 endpoints remain noisy (the body is mis-shaped at
  the cervicothoracic junction).
  \item \textbf{SPINEPS corpus border} --- fit to the SPINEPS corpus (label 49) border. This was a
  \emph{negative} result: noisier than canal-cut, because differencing two vertebrae amplifies a single
  bad per-vertebra fit.
  \item \textbf{SPINEPS endplate voxels (C1)} --- fit the line to SPINEPS' \emph{own predicted endplate
  voxels} (instance label $100+X$). This is the literature-validated object (the endplate itself, not the
  corpus border) and it won.
\end{enumerate}

\paragraph{The C1 method.} For each vertebra, the PCA major axis of the thin endplate sheet on the
mid-sagittal slice is the endplate tangent; the Cobb angle is the angle between two vertebrae's
inferior-endplate tangents, sign-calibrated to lordosis-positive. On the 12 healthy necks, C1 beats both
alternatives on every span: C2--C7 $\SI{+15.4}{\degree}\pm\SI{13.7}{\degree}$ (matching the F1000
literature mean of \SI{15.4}{\degree}), and the C6--C7 cervicothoracic endpoint --- the failure point of
the other methods --- drops to SD \SI{5.9}{\degree} (versus \SI{18.5}{\degree} for canal-cut). The
benchmark ceiling for an automated cervical-T2 Cobb is Zhang~et~al.\ (MAE \SI{2.44}{\degree})~\cite{zhang2025}.

\paragraph{Reporting and a corrected assumption.} Because absolute accuracy needs radiologist GT we lack,
Group 4 reports alignment as a descriptive class (lordotic/straightened/kyphotic) for review, with the
endpoint and supine caveats, rather than a hard-flagged magnitude. We also corrected a common false
assumption: supine MRI does \emph{not} read $\sim$\SI{5}{\degree} less lordotic than standing for C2--C7
(the difference is not significant in the F1000 cohort), so a flat supine$\rightarrow$standing offset
should not be applied.

\paragraph{Status.} C1 is the validated method (precision $+$ correct lordosis), confirmed on the
independent unhealthy cohort where it reads healthy $+15.2$ vs.\ unhealthy $+8.8$
(Section~\ref{sec:results}); as a \emph{discriminator} the separation is directional (loss of lordosis),
not significant at the current $n$, consistent with alignment being less specific for myelopathy than
stenosis.
